% =====================================================================
%  capitulos/01_introduccion.tex  --  SOLO CONTENIDO
% =====================================================================
%  Este fichero contiene únicamente texto. Lo que NO hay (a propósito):
%    - \documentclass, \usepackage ni \begin{document}: eso es de main.tex
%    - colores, márgenes ni formatos: eso es diseño y vive en el .sty
%  PISTAS PARA EDITAR
%    - Una línea en blanco separa párrafos.
%    - ~ es un espacio que NO permite salto de línea ("capítulo~2").
%    - \emph{...} = cursiva; \textbf{...} = negrita.
% =====================================================================
% \chapter crea un capítulo numerado y lo añade al índice.
\chapter{Introducción}

\section{Motivación}
La programación paralela es hoy una herramienta cotidiana, pero escribir
código paralelo eficiente y portable sigue siendo difícil. Una línea de
trabajo de la Universidad de Valladolid propone describir los programas
con especificaciones de alto nivel y dejar que un compilador modular
genere el código final~\cite{gonzalez2011trasgo}.
% \cite{clave} inserta una cita. La clave debe existir en bib/refs.bib; si
% no, saldrá [?]. Para varias a la vez: \cite{clave1,clave2}

% Entorno "aviso" definido en mistilo.sty (sustituye al \fcolorbox).
\begin{aviso}
  Esta memoria es un ejemplo para el taller: las referencias son reales;
  el resto, incluidos los datos, es ficticio.
\end{aviso}

\section{Objetivos}
El objetivo general es diseñar y evaluar un prototipo de consulta paralela.
Los objetivos específicos son:
\begin{enumerate}
  \item estudiar el estado del arte en paralelismo y datos enlazados;
  \item diseñar una arquitectura modular;
  \item implementar el prototipo en Python;
  \item medir su aceleración con distinto número de procesos.
\end{enumerate}

\section{Estructura de la memoria}
% \ref{etiqueta} se sustituye por el número de lo marcado con \label.
% Si se reordenan los capítulos, los números se actualizan solos.
El capítulo~\ref{cap:arte} revisa el estado del arte; el
capítulo~\ref{cap:diseno} presenta el diseño; el capítulo~\ref{cap:impl}
detalla la implementación; el capítulo~\ref{cap:result} muestra los
resultados y el capítulo~\ref{cap:conc} expone las conclusiones.
